\documentclass[11pt]{article}

\usepackage[a4paper,margin=30mm]{geometry}
\usepackage[T1]{fontenc}
\usepackage[utf8]{inputenc}
\usepackage{lmodern}
\usepackage{microtype}
\usepackage{amsmath,amssymb,amsthm,mathtools}
\usepackage{booktabs}
\usepackage{enumitem}
\usepackage{xcolor}
\usepackage[hidelinks]{hyperref}

\newtheorem{theorem}{Theorem}[section]
\newtheorem{proposition}[theorem]{Proposition}
\newtheorem{corollary}[theorem]{Corollary}
\newtheorem{lemma}[theorem]{Lemma}
\theoremstyle{definition}
\newtheorem{definition}[theorem]{Definition}
\newtheorem{assumption}[theorem]{Assumption}
\theoremstyle{remark}
\newtheorem{remark}[theorem]{Remark}
\newtheorem{question}[theorem]{Open question}

\DeclareMathOperator{\Fix}{Fix}
\DeclareMathOperator{\diag}{diag}
\newcommand{\R}{\mathbb{R}}
\newcommand{\N}{\mathcal{N}}
\newcommand{\pbar}{\bar p}
\newcommand{\one}{\mathbf{1}}
\newcommand{\pos}[1]{\left[#1\right]^+}

\title{Fractional Prepayment Games in Financial Networks\\
\large A reconstruction and extension of the 2020 LMFI notes}
\author{vd}
\date{2 September 2026}

\begin{document}

\maketitle

\begin{abstract}
This note develops the fractional-prepayment direction suggested by a set of
2020 working notes on clearing in interbank debt networks.  The underlying
model is the Rogers--Veraart extension of Eisenberg--Noe, with losses on
external assets and interbank receipts in default.  Before clearing, each bank
may use liquid external assets to repay any fraction of an existing liability.
The resulting game is a continuous counterpart of the all-or-nothing
prepayment game introduced in the later literature.

The first contribution is structural.  On every regular region with a fixed
default set, clearing payments are rational, and hence differentiable,
functions of selective prepayments.  They are affine when prepayments preserve
each debtor's liability proportions.  An explicit resolvent formula gives the
Jacobian.  The second contribution is local mechanism design.  For one
insolvent debtor, the note derives the exact cone of infinitesimal repayment
directions that weakly benefit every creditor; pro-rata repayment is the unique
direction that preserves all residual liability shares.  The third
contribution is strategic.  A four-node example has a unique best response that
prepays a strict fraction of one liability, showing that fractional strategies
cannot in general be reduced to all-or-nothing choices.  Earlier results from
the notes on fusion and round-trip loans are retained as background.
\end{abstract}

\tableofcontents

\section{Purpose, contribution, and related literature}

The original notes move between a clearing model, comparative statics, and a
game in which banks alter their balance sheets before clearing.  Several
calculations are exploratory, and some statements labelled as theorems are
followed by requests to verify them.  The aim here is to make the assumptions
visible and to pursue the strongest undeveloped direction: continuous
prepayment before clearing.

The main conclusions are as follows.

\begin{enumerate}[label=(\roman*)]
  \item The greatest clearing fixed point is the natural selection when bank
  value is increasing in receipts from the network.

  \item Fusion has a transparent local accounting rule.  The deficit of the
  fused node is the sum of the two original deficits and the fusion cost,
  minus the two unpaid internal claims that disappear in the fusion.

  \item In the restricted three-node model, balance-sheet-consistent
  round-trip lending cannot increase the solvent accomplice's payoff.  The
  conclusion in the thesis is valid, but the asymmetric payoff formula needs
  a cash-flow correction.

  \item Selective early repayment creates aggregate value when liquidation of
  external assets is costly.  The selected creditor weakly gains, but other
  creditors may lose.  The exact local Pareto cone describes how much must be
  repaid to each creditor to eliminate this conflict.

  \item A continuous prepayment game has compact polyhedral strategy sets.
  Within a regular fixed-default region its payoffs are differentiable rational
  functions.  A resolvent formula reduces every local marginal calculation to
  one linear solve.

  \item Fractional actions are essential.  There are networks in which a
  player's unique best response repays neither zero nor the whole liability.
  The optimum is the least payment that moves another bank across a solvency
  boundary.
\end{enumerate}

Statements about derivatives are local to a region in which the same banks
remain solvent or insolvent.  Crossing a default boundary changes the formula
for the clearing map and must be handled separately.

\subsection{Relation to prior work}

Eisenberg and Noe established the fixed-point clearing model with
proportionality, limited liability, and priority of debt
\cite{EisenbergNoe2001}.  Rogers and Veraart introduced default costs and
studied mergers and rescue consortia \cite{RogersVeraart2013}.  The fusion
material below therefore sharpens an existing operation rather than proposing
a new one.  Netting and portfolio compression had also shown before 2020 that
removing gross or circular liabilities can have adverse distributional and
systemic effects \cite{AminiEtAl2016,SchuldenzuckerSeuken2020}.

Strategic departures from proportional clearing were studied by Bertschinger,
Hoefer, and Schmand \cite{BertschingerEtAl2020}.  Papp and Wattenhofer showed
that debt forgiveness and transfers of external assets can improve a bank's
payoff through network feedback \cite{PappWattenhofer2020}.  These operations
are close to prepayment, but they do not model the same balance-sheet update at
an earlier date.  Work on multiple maturities supplies a dynamic clearing
background \cite{KusnetsovVeraart2019}.

Sensitivity formulas for Eisenberg--Noe payments predate the present note.
Liu and Staum differentiate with respect to external assets
\cite{LiuStaum2010}; Feinstein et al. differentiate with respect to bilateral
liabilities \cite{FeinsteinEtAl2018}.  The derivative formula below applies
the same general idea to the simultaneous change in assets, nominal debts, and
relative liabilities caused by prepayment.

Zhou et al. later introduced a strategic prepayment game based on
all-or-nothing repayment of whole liabilities \cite{ZhouEtAl2024}.  Their
conclusion lists fractional prepayments and coalitional incentives as open
directions.  The 2020 notes appear to anticipate this topic, but the present
paper makes no claim of publication priority.  Its narrower objective is to
develop the continuous model and results that are absent from the discrete
formulation.

\section{Clearing with default costs}

\subsection{Debt networks}

Let the set of banks be $\N=\{1,\ldots,n\}$.  A debt network consists of
external assets $e\in\R_+^n$ and a liability matrix
$L\in\R_+^{n\times n}$ with $L_{ii}=0$.  The entry $L_{ij}$ is the nominal
amount owed by bank $i$ to bank $j$.  Define total nominal liabilities and
relative liabilities by
\begin{equation}
  \pbar_i=\sum_{j\in\N}L_{ij},
  \qquad
  \pi_{ij}=
  \begin{cases}
    L_{ij}/\pbar_i, & \pbar_i>0,\\
    0, & \pbar_i=0.
  \end{cases}
\end{equation}
Thus $\pi_{ij}p_i$ is the payment from $i$ to $j$ when $i$ makes the
total payment $p_i$.

Let $\alpha,\beta\in[0,1]$.  The parameter $\alpha$ is the recovery fraction
on external assets in default, and $\beta$ is the recovery fraction on
interbank receipts.  The Rogers--Veraart clearing map
$\Phi:[0,\pbar]\to[0,\pbar]$ is
\begin{equation}\label{eq:clearing-map}
  \Phi_i(p)=
  \begin{cases}
    \pbar_i,
      & e_i+\displaystyle\sum_{j\in\N}\pi_{ji}p_j\geq\pbar_i,\\[2mm]
    \alpha e_i+\beta\displaystyle\sum_{j\in\N}\pi_{ji}p_j,
      & e_i+\displaystyle\sum_{j\in\N}\pi_{ji}p_j<\pbar_i.
  \end{cases}
\end{equation}
A clearing vector is a fixed point of $\Phi$.

The map is monotone in the product order.  Tarski's fixed-point theorem gives
a least and a greatest fixed point.  We denote the greatest one by $p^\uparrow$.
It can also be obtained by the decreasing iteration or default-set algorithm
described by Rogers and Veraart \cite{RogersVeraart2013}.

\subsection{Why select the greatest fixed point?}

For a payment vector $p$, define the post-clearing equity of bank $i$ by
\begin{equation}\label{eq:equity}
  V_i(p)=
  \pos{e_i+\sum_{j\in\N}\pi_{ji}p_j-\pbar_i}.
\end{equation}
This definition treats debt payments as senior to equity.

\begin{proposition}[Pareto dominance of the greatest fixed point]
For every clearing vector $p$, one has
\begin{equation}
  V_i(p^\uparrow)\geq V_i(p)
  \qquad\text{for every }i\in\N.
\end{equation}
\end{proposition}

\begin{proof}
By definition of the greatest fixed point, $p^\uparrow\geq p$ coordinatewise.
Every coefficient $\pi_{ji}$ is nonnegative, so the expression inside the
positive part in \eqref{eq:equity} is weakly increasing in $p$.  The positive
part is also increasing.
\end{proof}

The selection matters when default costs destroy uniqueness.  For example,
take two banks with
\begin{equation}
  L_{12}=L_{21}=2.2,
  \qquad e_1=e_2=1,
  \qquad \alpha=\beta=0.5.
\end{equation}
Both $(1,1)$ and $(2.2,2.2)$ are clearing vectors.  The second is the greatest
fixed point.  This example also shows why an iteration from zero cannot simply
be interpreted as selecting the economically preferred outcome.

\section{Fusion and the cancellation of internal claims}

\subsection{Definition of fusion}

Fix two distinct banks $a,b\in\N$.  Their fusion is a new node $[ab]$.  Mutual
liabilities $L_{ab}$ and $L_{ba}$ disappear, while liabilities to and from all
other banks are added.  If $\kappa\geq0$ is the total fusion cost, then
\begin{align}
  e_{[ab]}&=e_a+e_b-\kappa,\\
  \pbar_{[ab]}&=\pbar_a+\pbar_b-L_{ab}-L_{ba},\\
  L_{[ab]j}&=L_{aj}+L_{bj},\qquad j\notin\{a,b\},\\
  L_{j[ab]}&=L_{ja}+L_{jb},\qquad j\notin\{a,b\}.
\end{align}

The full clearing vector may change after fusion because the fused node has a
new liability proportion.  It is therefore useful to begin with a local
identity at fixed external payments.

\subsection{A fusion deficit identity}

For any candidate payment vector $p\in[0,\pbar]$, define the pre-haircut
resource deficit of bank $i$ by
\begin{equation}
  \delta_i(p)=
  \pbar_i-e_i-\sum_{j\in\N}\pi_{ji}p_j.
\end{equation}
Positive $\delta_i(p)$ means that $i$ cannot pay in full at the candidate
payments.  Define the unpaid portions of the two internal claims by
\begin{equation}
  u_{ab}(p)=L_{ab}-\pi_{ab}p_a,
  \qquad
  u_{ba}(p)=L_{ba}-\pi_{ba}p_b.
\end{equation}
Both quantities are nonnegative.

\begin{theorem}[Fusion deficit identity]\label{thm:fusion-identity}
Hold the payments of all banks outside $\{a,b\}$ fixed.  The deficit of the
fused node is
\begin{equation}\label{eq:fusion-identity}
  \delta_{[ab]}
  =\delta_a(p)+\delta_b(p)+\kappa-u_{ab}(p)-u_{ba}(p).
\end{equation}
\end{theorem}

\begin{proof}
The external inflow of the fused node is
\begin{equation}
  \sum_{j\notin\{a,b\}}p_j(\pi_{ja}+\pi_{jb}).
\end{equation}
Consequently,
\begin{align*}
  \delta_{[ab]}
  ={}&\pbar_a+\pbar_b-L_{ab}-L_{ba}
      -(e_a+e_b-\kappa)\\
    &-\sum_{j\notin\{a,b\}}p_j(\pi_{ja}+\pi_{jb}).
\end{align*}
On the other hand,
\begin{align*}
  \delta_a(p)+\delta_b(p)
  ={}&\pbar_a+\pbar_b-e_a-e_b\\
    &-\sum_{j\notin\{a,b\}}p_j(\pi_{ja}+\pi_{jb})
      -\pi_{ab}p_a-\pi_{ba}p_b.
\end{align*}
Subtracting these expressions gives
\begin{align*}
  \delta_{[ab]}-\delta_a(p)-\delta_b(p)
   & =\kappa-L_{ab}-L_{ba}+\pi_{ab}p_a+\pi_{ba}p_b\\
   & =\kappa-u_{ab}(p)-u_{ba}(p).
\end{align*}
\end{proof}

The identity makes the economics explicit.  Fusion adds its direct cost, but
it also deletes internal claims that were not going to be paid.  Those unpaid
claims are the only local source of a rescue benefit.

\begin{corollary}[Local rescue criterion]\label{cor:fusion-rescue}
At fixed external payments, the fused node can pay in full if and only if
\begin{equation}
  u_{ab}(p)+u_{ba}(p)
  \geq \delta_a(p)+\delta_b(p)+\kappa.
\end{equation}
\end{corollary}

\begin{corollary}[Common-stage insolvency]
Suppose $a$ and $b$ are still nominal payers at a given stage, so that
$p_a=\pbar_a$ and $p_b=\pbar_b$.  If
$\delta_a(p)>0$ and $\delta_b(p)>0$, then their fusion is insolvent at that
stage for every $\kappa\geq0$.
\end{corollary}

\begin{proof}
Nominal payment implies $u_{ab}(p)=u_{ba}(p)=0$.  Apply
Theorem~\ref{thm:fusion-identity}.
\end{proof}

This corollary gives a precise version of the claim that two banks becoming
insolvent at the same default-set iteration cannot rescue each other by
fusion.  A global theorem additionally requires checking that payments of the
other banks agree up to that iteration.  This agreement holds under the usual
same-stage construction before either fused bank reduces its payment, but it
should not be assumed after their default because the fused liability shares
then change.

\begin{question}[Global fusion comparative statics]
Can one characterize the change in the greatest clearing vector after fusion
using the resolvent of the default-set linear system?  The identity
\eqref{eq:fusion-identity} gives the direct effect.  A complete answer must also
capture feedback through the payments of all other defaulting banks.
\end{question}

\section{Round-trip loans in a three-node network}

\subsection{The restricted model}

Consider banks $1$, $2$, and $n$.  Bank $1$ owes $B>0$ to bank $2$ and $A\geq0$
to bank $n$.  Bank $2$ owes $C\geq0$ to bank $1$.  Bank $1$ has no initial
external assets.  Banks $2$ and $n$ remain solvent, while bank $1$ remains
insolvent.  The last condition is
\begin{equation}\label{eq:base-insolvency}
  C<A+B.
\end{equation}

Suppose bank $n$ lends $S_1\geq0$ to bank $1$, and bank $1$ lends
$S_2\in[0,S_1]$ back to bank $n$.  Balance-sheet consistency gives bank $1$
the external assets $S_1-S_2$.  Its total payment in default is
\begin{equation}
  \alpha(S_1-S_2)+\beta(C+S_2),
\end{equation}
and the fraction directed to bank $n$ is
\begin{equation}
  q(S_1)=\frac{A+S_1}{A+B+S_1}.
\end{equation}

The relevant net cash contribution to bank $n$ is therefore
\begin{equation}\label{eq:corrected-loan-payoff}
  \nu(S_1,S_2)
  =\bigl[\alpha(S_1-S_2)+\beta(C+S_2)\bigr]q(S_1)-S_1.
\end{equation}
Only $S_1$ is subtracted.  The cash $S_2$ borrowed by bank $n$ is offset by
bank $n$'s repayment of the same debt.  Subtracting both $S_1$ and $S_2$, as in
the exploratory formula, counts the second flow twice.

\subsection{No profitable round trip}

\begin{theorem}[Round-trip lending is unprofitable]\label{thm:round-trip}
Under \eqref{eq:base-insolvency}, for every
$\alpha,\beta\in[0,1]$ and every feasible pair
$0\leq S_2\leq S_1$,
\begin{equation}
  \nu(S_1,S_2)\leq\nu(0,0).
\end{equation}
\end{theorem}

\begin{proof}
For fixed $S_1$, the function $\nu$ is affine in $S_2$, with
\begin{equation}
  \frac{\partial\nu}{\partial S_2}
  =(\beta-\alpha)q(S_1).
\end{equation}

First suppose $\beta\geq\alpha$.  The maximum over $S_2\in[0,S_1]$ is at
$S_2=S_1$.  Put $S=S_1$ and
\begin{equation}
  f(S)=\beta(S+C)\frac{A+S}{A+B+S}-S.
\end{equation}
The derivative of the fraction before multiplication by $\beta$ is positive,
so $\beta\leq1$ gives
\begin{equation}
  f'(S)
  \leq
  \frac{B(C-A-B)}{(A+B+S)^2}<0.
\end{equation}

Now suppose $\beta<\alpha$.  The maximum over $S_2$ is at $S_2=0$.  Put
\begin{equation}
  g(S)=\bigl(\alpha S+\beta C\bigr)
       \frac{A+S}{A+B+S}-S.
\end{equation}
All coefficients multiplying $\alpha$ and $\beta$ in $g'(S)$ are
nonnegative.  Replacing both parameters by $1$ gives an upper bound, hence
\begin{equation}
  g'(S)
  \leq
  \frac{B(C-A-B)}{(A+B+S)^2}<0.
\end{equation}
In both cases the best value is attained at $S_1=S_2=0$.
\end{proof}

Theorem~\ref{thm:round-trip} includes the symmetric case $S_1=S_2$.  Its scope
is intentionally narrow.  A larger network may change the solvency of bank
$2$ or create feedback through other defaulting nodes.  Those regime changes
are not represented by \eqref{eq:corrected-loan-payoff}.

\begin{remark}[Admissible balance-sheet changes]
Changing entries of $L$ without the corresponding transfer of assets creates
nominal claims from nothing.  Such an operation can alter proportional shares
and can look profitable, but it is not a loan.  Any strategic model should
require double-entry balance-sheet consistency for every action.
\end{remark}

\section{Selective early repayment}

\subsection{Fixed-regime setup}

Let bank $0$ be insolvent and let its creditors be $1,\ldots,m$.  Write
\begin{equation}
  P=\sum_{i=1}^m L_{0i}
\end{equation}
for its initial liabilities.  Before clearing, bank $0$ makes early repayments
$x_i$ satisfying
\begin{equation}\label{eq:repayment-feasibility}
  0\leq x_i\leq L_{0i},
  \qquad
  X:=\sum_{i=1}^m x_i\leq e_0.
\end{equation}
Let $I$ denote the fixed incoming payment to bank $0$ from the rest of the
network.  Assume throughout this section that the creditors remain solvent and
that the other banks' payments do not change as $x$ varies.

After repayment, bank $0$ has
\begin{equation}
  E=e_0-X,
  \qquad
  D=P-X,
\end{equation}
where $E$ is its remaining external asset and $D$ its remaining nominal debt.
If bank $0$ was initially 0-insolvent, then
\begin{equation}
  e_0+I<P.
\end{equation}
Subtracting $X$ from both sides gives $E+I<D$.  Early repayment cannot make
bank $0$ solvent by itself.

The total clearing payment made by bank $0$ on its remaining debt is
\begin{equation}
  W=\alpha E+\beta I.
\end{equation}
Since $\alpha,\beta\leq1$, one has $W<D$.  Define the remaining liability share
and effective recovery rate by
\begin{equation}
  q_i=\frac{L_{0i}-x_i}{D},
  \qquad
  r=\frac{W}{D}.
\end{equation}
Then $q_i\in[0,1]$, $\sum_iq_i=1$, and $r<1$.

Up to terms independent of $x$, creditor $i$ has value
\begin{equation}\label{eq:creditor-value}
  R_i(x)=x_i+Wq_i+K_i,
\end{equation}
where $K_i$ collects its external assets, other incoming payments, and
liabilities.

\subsection{Own and cross effects}

\begin{theorem}[Marginal effects of selective repayment]
Within the fixed default regime,
\begin{align}
  \frac{\partial R_i}{\partial x_i}
    &=1-\alpha q_i-r(1-q_i)\geq0,
      \label{eq:own-effect}\\
  \frac{\partial R_i}{\partial x_k}
    &=q_i(r-\alpha),
      \qquad k\neq i,
      \label{eq:cross-effect}\\
  \frac{\partial}{\partial x_k}\sum_{i=1}^m R_i
    &=1-\alpha\geq0.
      \label{eq:aggregate-effect}
\end{align}
\end{theorem}

\begin{proof}
For repayment to creditor $i$, both $L_{0i}-x_i$ and $D$ have derivative
$-1$, while $W$ has derivative $-\alpha$.  Differentiating
\eqref{eq:creditor-value} gives
\begin{align*}
  \frac{\partial R_i}{\partial x_i}
  &=1-\alpha\frac{L_{0i}-x_i}{D}
    -W\frac{D-(L_{0i}-x_i)}{D^2}\\
  &=1-\alpha q_i-r(1-q_i).
\end{align*}
The quantity subtracted from $1$ is a convex combination of $\alpha$ and
$r$, both of which are at most $1$.  This proves
\eqref{eq:own-effect}.

For $k\neq i$, the numerator $L_{0i}-x_i$ is fixed, while $W$ and $D$ both
decrease.  Hence
\begin{equation*}
  \frac{\partial R_i}{\partial x_k}
  =-\alpha q_i+\frac{W}{D}q_i=q_i(r-\alpha),
\end{equation*}
which is \eqref{eq:cross-effect}.  Finally,
\begin{equation}
  \sum_{i=1}^m R_i=X+W+\sum_{i=1}^m K_i,
\end{equation}
and differentiation gives \eqref{eq:aggregate-effect}.
\end{proof}

Equation~\eqref{eq:own-effect} proves the early-repayment conjecture in the
2020 thesis under its stated fixed-regime assumptions.  The selected creditor
cannot be harmed by receiving more money early.

Equation~\eqref{eq:cross-effect} gives the missing distributional result.
All creditors not selected for repayment move in the same direction:
\begin{itemize}
  \item if $r<\alpha$, every non-selected creditor with a positive remaining
  claim loses;
  \item if $r=\alpha$, non-selected creditors are locally indifferent;
  \item if $r>\alpha$, every creditor weakly benefits.
\end{itemize}

\begin{corollary}[Marginal Pareto criterion]
A selective early repayment is a Pareto improvement for the creditors if
$r\geq\alpha$.  If $r<\alpha$ and a non-selected creditor has a positive
remaining claim, the repayment is not a Pareto improvement even though total
creditor value rises by $1-\alpha$ per unit.
\end{corollary}

The conflict is therefore not between efficiency and all repayment.  It is
between aggregate recovery and the distribution of that recovery.

The preceding calculation extends from a payment to one creditor to an
arbitrary infinitesimal repayment direction.  Let $h_i\geq0$ be the rate of
repayment to creditor $i$ and put $H=\sum_i h_i$.

\begin{theorem}[The local Pareto cone]\label{thm:pareto-cone}
Within the fixed default regime, the directional derivative of creditor $i$'s
value is
\begin{equation}\label{eq:directional-creditor}
  \dot R_i=(1-r)h_i+(r-\alpha)Hq_i.
\end{equation}
Consequently:
\begin{enumerate}[label=(\roman*)]
  \item if $r\geq\alpha$, every nonnegative repayment direction is locally
  Pareto improving for the creditors;
  \item if $r<\alpha$ and $H>0$, a direction is locally Pareto improving if
  and only if
  \begin{equation}\label{eq:pareto-cone}
    \frac{h_i}{H}
    \geq q_i\frac{\alpha-r}{1-r}
    \qquad\text{for every }i.
  \end{equation}
\end{enumerate}
The cone in \eqref{eq:pareto-cone} is nonempty.  When $\alpha<1$, the share of
the marginal repayment not forced by the inequalities is
$(1-\alpha)/(1-r)$.
\end{theorem}

\begin{proof}
Along the direction $h$, one has
\begin{equation}
  \dot W=-\alpha H,
  \qquad
  \dot D=-H,
  \qquad
  \dot q_i=\frac{Hq_i-h_i}{D}.
\end{equation}
Differentiating $R_i=x_i+Wq_i+K_i$ gives
\begin{equation*}
  \dot R_i
  =h_i-\alpha Hq_i+r(Hq_i-h_i),
\end{equation*}
which is \eqref{eq:directional-creditor}.  If $r\geq\alpha$, both terms in
\eqref{eq:directional-creditor} are nonnegative.  If $r<\alpha$, then $r<1$,
and rearranging $\dot R_i\geq0$ gives \eqref{eq:pareto-cone}.  The lower bounds
on the repayment shares sum to
\begin{equation}
  \frac{\alpha-r}{1-r}\sum_iq_i
  =\frac{\alpha-r}{1-r}\leq1.
\end{equation}
The remaining share equals $(1-\alpha)/(1-r)$.
\end{proof}

Theorem~\ref{thm:pareto-cone} gives an exact local compensation rule.  When
$r<\alpha$, a targeted repayment may still be safe, but excluded creditors
must receive the lower-bound shares in \eqref{eq:pareto-cone}.  Any remaining
share can be directed to the creditor whose rescue creates the largest network
benefit.

\subsection{A pro-rata safe harbor}

Selective repayment changes liability shares.  A natural alternative is to
allow early repayment only in the original proportions.

\begin{theorem}[Pro-rata early repayment]\label{thm:pro-rata}
Let
\begin{equation}
  x_i(t)=tL_{0i},
  \qquad 0\leq t<1,
  \qquad tP\leq e_0.
\end{equation}
Within the fixed default regime,
\begin{equation}\label{eq:pro-rata-gain}
  \frac{dR_i}{dt}=(1-\alpha)L_{0i}\geq0
  \qquad\text{for every creditor }i.
\end{equation}
Thus pro-rata early repayment weakly benefits every creditor and preserves
their relative claims.
\end{theorem}

\begin{proof}
The remaining debt is $D(t)=(1-t)P$, and
\begin{equation}
  q_i(t)=\frac{(1-t)L_{0i}}{(1-t)P}=\frac{L_{0i}}{P}
\end{equation}
is constant.  Also
\begin{equation}
  W(t)=\alpha(e_0-tP)+\beta I.
\end{equation}
Substitution in \eqref{eq:creditor-value} gives
\begin{equation}
  R_i(t)=tL_{0i}+\frac{L_{0i}}{P}W(t)+K_i.
\end{equation}
Differentiating proves \eqref{eq:pro-rata-gain}.
\end{proof}

\begin{proposition}[Characterization by share preservation]
Among nonzero infinitesimal repayment directions, the pro-rata direction
$h_i=Hq_i$ is the unique direction that leaves every residual liability share
$q_i$ unchanged.
\end{proposition}

\begin{proof}
The derivative calculated in the proof of
Theorem~\ref{thm:pareto-cone} satisfies
\begin{equation}
  \dot q_i=0
  \quad\Longleftrightarrow\quad
  h_i=Hq_i.
\end{equation}
Applying this equivalence to every $i$ proves the claim.
\end{proof}

Theorem~\ref{thm:pro-rata} suggests a simple rule: a standstill need not ban
all early repayment.  It can permit payments that preserve the original
creditor proportions.  Such payments capture the value otherwise lost through
the factor $\alpha$ without preferring one creditor over another.

\subsection{Rescuing a creditor exposed to contagion}

Suppose creditor $i$ would become insolvent because bank $0$ defaults.  With
all other early repayments fixed, Equation~\eqref{eq:own-effect} shows that
$R_i$ is nondecreasing in $x_i$.  Consequently the exact rescue test is an
endpoint calculation:
\begin{equation}
  R_i\bigl(x_i^{\max}\bigr)\geq0,
  \qquad
  x_i^{\max}=\min\left\{L_{0i},
  e_0-\sum_{k\neq i}x_k\right\}.
\end{equation}
If the inequality fails, no feasible selective repayment to $i$ can rescue it
within the fixed regime.  If it holds, continuity gives a least rescue
threshold.  This formulation is more precise than a sufficient inequality
that ignores the creditor's other assets and liabilities.

\section{The continuous prepayment game}\label{sec:continuous-game}

\subsection{Strategies and balance-sheet updates}

Fix an initial debt network $(e,L)$.  Bank $i$ chooses a vector of outgoing
prepayments
\begin{equation}
  x_i=(x_{ij})_{j\in\N}
\end{equation}
from the compact polytope
\begin{equation}\label{eq:strategy-polytope}
  \mathcal X_i=
  \left\{x_i:
  0\leq x_{ij}\leq L_{ij},\quad
  \sum_jx_{ij}\leq b_i\right\},
  \qquad 0\leq b_i\leq e_i.
\end{equation}
Using an initial liquidity budget $b_i$ makes the joint strategy space the
product $\mathcal X=\prod_i\mathcal X_i$.  In particular, one bank cannot fund
its simultaneous strategy using a prepayment that it receives from another
bank in the same action round.

For a profile $x\in\mathcal X$, double-entry consistency gives
\begin{align}
  e_i^x
    &=e_i-\sum_jx_{ij}+\sum_jx_{ji},
      \label{eq:updated-assets}\\
  L_{ij}^x
    &=L_{ij}-x_{ij},
      \label{eq:updated-liabilities}\\
  \pbar_i^x
    &=\sum_jL_{ij}^x.
      \label{eq:updated-total-liabilities}
\end{align}
Thus $\sum_i e_i^x=\sum_i e_i$: prepayment transfers external assets but does
not create them.  Let $\Pi^x$ be the relative-liability matrix constructed
from $L^x$, and let $p^\uparrow(x)$ be the greatest Rogers--Veraart clearing
vector of the updated network.

Two natural utility functions are total assets and equity:
\begin{align}
  A_i(x)
    &=e_i^x+\sum_j\pi_{ji}^x p_j^\uparrow(x),
      \label{eq:total-asset-utility}\\
  V_i(x)
    &=\pos{A_i(x)-\pbar_i^x}.
      \label{eq:equity-utility}
\end{align}
Total assets distinguish outcomes for an insolvent bank whose equity is zero.
Equity is the direct objective of shareholders.  Since these objectives can
produce different games, every equilibrium statement should specify which one
is being used.

\subsection{Regular default regions}

Let $D$ be a proposed set of defaulting banks and let $S=\N\setminus D$.
Consider an open set of prepayment profiles on which all banks in $D$ are
strictly insolvent, all banks in $S$ are solvent, and
$\pbar_i^x>0$ for every relevant row.  On this set,
\begin{equation}\label{eq:default-cell-system}
  p_D(x)
  =\alpha e_D^x
   +\beta(\Pi_{DD}^x)^{\mathsf T}p_D(x)
   +\beta(\Pi_{SD}^x)^{\mathsf T}\pbar_S^x,
  \qquad
  p_S(x)=\pbar_S^x.
\end{equation}
Define the default-set resolvent
\begin{equation}\label{eq:default-resolvent}
  M_D(x)=
  \left[I-\beta(\Pi_{DD}^x)^{\mathsf T}\right]^{-1}
\end{equation}
whenever the inverse exists.  It always exists when $\beta<1$ because
$\Pi_{DD}^x$ is substochastic.  When $\beta=1$, invertibility is a regularity
condition on the defaulting subnetwork.

\begin{theorem}[Geometry of a fixed-default region]
\label{thm:fixed-default-geometry}
On a region satisfying the preceding strict inequalities and invertibility
condition,
\begin{equation}\label{eq:fixed-cell-solution}
  p_D(x)=M_D(x)
  \left[
    \alpha e_D^x
    +\beta(\Pi_{SD}^x)^{\mathsf T}\pbar_S^x
  \right].
\end{equation}
The clearing payments and the utilities in
\eqref{eq:total-asset-utility}--\eqref{eq:equity-utility} are rational, hence
real analytic, functions of $x$ on the region.  They need not be affine.

If every active debtor prepays proportionally,
\begin{equation}\label{eq:row-proportional-action}
  x_{ij}=t_iL_{ij},
  \qquad 0\leq t_i<1,
\end{equation}
then $\Pi^x=\Pi$.  On every fixed-default region, $p^\uparrow(x)$ is then an
affine function of the vector $t$.
\end{theorem}

\begin{proof}
Equation~\eqref{eq:fixed-cell-solution} follows by rearranging
\eqref{eq:default-cell-system}.  Assets and nominal liabilities are affine in
$x$.  Every nonzero relative-liability entry has the form
\begin{equation}
  \pi_{ij}^x=\frac{L_{ij}-x_{ij}}
  {\pbar_i-\sum_kx_{ik}},
\end{equation}
so it is rational.  Matrix inversion expresses each entry of $M_D(x)$ as a
ratio of determinants.  Formula~\eqref{eq:fixed-cell-solution} is therefore
rational wherever the denominators and determinant do not vanish.  Selective
prepayment generally changes the displayed ratio nonlinearly, which rules out
a general affine claim.

Under \eqref{eq:row-proportional-action}, every residual row of liabilities is
multiplied by $1-t_i$, so its relative proportions are unchanged.  The matrix
in \eqref{eq:default-resolvent} is constant, while $e^x$ and $\pbar^x$ are
affine in $t$.  Formula~\eqref{eq:fixed-cell-solution} is therefore affine.
\end{proof}

Theorem~\ref{thm:fixed-default-geometry} gives a finite default-set
decomposition, but it is not a globally smooth decomposition.  Default
costs can create jumps when a bank becomes just solvent.  Such boundaries can
determine best responses, as Theorem~\ref{thm:fractional-best-response} below
shows.

\subsection{A resolvent formula for marginal effects}

Let $x(t)=x+th$ remain inside a regular fixed-default region for small $t$.
Put
\begin{equation}
  H_i=\sum_jh_{ij}.
\end{equation}
Dots denote derivatives at $t=0$.  Direct differentiation of the
balance-sheet update gives
\begin{align}
  \dot e_i&=-H_i+\sum_jh_{ji},
      \label{eq:asset-direction}\\
  \dot{\pbar}_i&=-H_i,
      \label{eq:liability-direction}\\
  \dot\pi_{ij}
    &=\frac{\pi_{ij}^xH_i-h_{ij}}{\pbar_i^x}.
      \label{eq:relative-liability-direction}
\end{align}

\begin{theorem}[Prepayment Jacobian]\label{thm:prepayment-jacobian}
Inside a regular region with default set $D$, the directional derivative of
the clearing vector is $\dot p_S=\dot{\pbar}_S$ and
\begin{equation}\label{eq:prepayment-jacobian}
  \dot p_D=M_D(x)\left[
    \alpha\dot e_D
    +\beta(\dot\Pi_{DD})^{\mathsf T}p_D
    +\beta(\dot\Pi_{SD})^{\mathsf T}\pbar_S
    +\beta(\Pi_{SD}^x)^{\mathsf T}\dot{\pbar}_S
  \right].
\end{equation}
Consequently,
\begin{equation}\label{eq:asset-utility-jacobian}
  \dot A_i
  =\dot e_i+
   \sum_j\left(\dot\pi_{ji}p_j+\pi_{ji}^x\dot p_j\right).
\end{equation}
\end{theorem}

\begin{proof}
Differentiate \eqref{eq:default-cell-system} and collect the terms containing
$\dot p_D$ on the left.  Multiplication by the inverse in
\eqref{eq:default-resolvent} gives
\eqref{eq:prepayment-jacobian}.  Equation
\eqref{eq:asset-utility-jacobian} follows by differentiating
\eqref{eq:total-asset-utility}.
\end{proof}

The resolvent $M_D$ is a network multiplier.  The bracketed vector in
\eqref{eq:prepayment-jacobian} is the direct balance-sheet and liability-share
effect of the action; multiplication by $M_D$ propagates it through cycles of
defaulting banks.  A local best-response calculation therefore requires one
factorization of $I-\beta(\Pi_{DD}^x)^{\mathsf T}$ and one linear solve per
candidate direction.

\subsection{A uniquely optimal fractional prepayment}

The discrete prepayment model permits only full repayment of an edge or no
repayment.  The following example shows that this restriction can exclude the
unique best response.

\begin{theorem}[Fractional best responses are essential]
\label{thm:fractional-best-response}
For every $\alpha\in[0,1)$, there is a Rogers--Veraart network with $\beta=0$
in which a bank maximizing total assets has a unique best response that
prepays a strict fraction of one liability.  The same conclusion holds for
every $\beta\in[0,1)$ when $\alpha=0$.
\end{theorem}

\begin{proof}
Consider banks $0,1,2$ and an outside sink $\omega$.  Let
\begin{equation}
  L_{01}=10,
  \qquad L_{0\omega}=30,
  \qquad L_{12}=10,
  \qquad L_{20}=10,
\end{equation}
with all other liabilities zero.  External assets are
\begin{equation}
  e_0=10,
  \qquad e_1=4,
  \qquad e_2=e_\omega=0,
\end{equation}
and first set $\beta=0$.  Only bank $0$ acts.  It chooses a prepayment
$x=x_{01}\in[0,10]$.  Bank $0$ is always insolvent and pays
$p_0=\alpha(10-x)$.  Its clearing payment to bank $1$ is
\begin{equation}
  z_{01}(x)
  =\alpha\frac{(10-x)^2}{40-x}.
\end{equation}
Bank $1$ becomes solvent at the unique solution $x_\alpha$ of
\begin{equation}\label{eq:fractional-threshold}
  4+x+\alpha\frac{(10-x)^2}{40-x}=10.
\end{equation}
Indeed, the left-hand side is strictly increasing on $[0,10]$.  To see this,
put $g(x)=(10-x)^2/(40-x)$.  Then
\begin{equation}
  g'(x)=\frac{(10-x)(x-70)}{(40-x)^2}
  \geq-\frac{7}{16},
\end{equation}
so the derivative of the left-hand side of
\eqref{eq:fractional-threshold} is at least $1-7\alpha/16>0$.  Its value at
$x=0$ is $4+5\alpha/2<10$, while its value at $x=6$ is
$10+8\alpha/17\geq10$.  Hence
\begin{equation}
  0<x_\alpha\leq6<10.
\end{equation}

Below $x_\alpha$, bank $1$ defaults and bank $2$ receives less than its
nominal claim.  Since $\beta=0$ and bank $2$ has no external assets, it pays
zero to bank $0$.  At and above $x_\alpha$, bank $1$ pays $10$ to bank $2$,
which becomes solvent and pays $10$ to bank $0$.  Bank $0$ remains insolvent
because its post-action liabilities are $40-x$ and its pre-haircut resources
are at most $20-x$.  Bank $0$'s total-asset utility is therefore
\begin{equation}
  A_0(x)=
  \begin{cases}
    10-x, & 0\leq x<x_\alpha,\\
    20-x, & x_\alpha\leq x\leq10.
  \end{cases}
\end{equation}
Its unique maximizer is $x=x_\alpha$, a strict fraction of $L_{01}$.

Finally, set $\alpha=0$ and take any $\beta<1$.  For $x<6$, the payments in
the cycle satisfy
\begin{equation}
  p_0=\beta p_2,
  \qquad
  p_1=\beta\frac{10-x}{40-x}p_0,
  \qquad
  p_2=\beta p_1.
\end{equation}
Their only solution is zero because
$\beta^3(10-x)/(40-x)<1$.  At $x=6$, bank $1$ and then bank $2$ become
solvent.  The same utility calculation gives a unique maximizer $x=6$.
\end{proof}

The optimum in this example is a solvency threshold.  Below it, the marginal
payment does not activate the return path $1\to2\to0$; above it, further
prepayment only spends bank $0$'s external assets.  This phenomenon is missed
both by endpoint-only strategies and by optimization confined to the interior
of a fixed-default region.

\section{A precise standstill problem}

The word ``fraud'' in the working notes covers two different phenomena.
Creating a liability without transferring the corresponding asset violates
balance-sheet consistency.  Selective early repayment is different: it is a
real transfer, but it violates equal treatment of creditors.

The latter creates a simple coalition incentive.

\begin{proposition}[Bilateral repayment incentive]
Suppose bank $0$ remains insolvent, so its post-clearing equity is zero.  If
bank $0$ and creditor $i$ can jointly choose $x_i$ at no transaction cost while
other repayments are fixed, then the coalition weakly prefers the largest
feasible $x_i$.  The preference is strict whenever creditor $i$'s derivative
in \eqref{eq:own-effect} is positive.
\end{proposition}

\begin{proof}
Bank $0$ remains 0-insolvent because early repayment subtracts the same amount
from its assets and liabilities.  Its equity therefore remains zero.  The
selected creditor's value is nondecreasing by
\eqref{eq:own-effect}.
\end{proof}

This proposition captures the valid part of the tentative standstill theorem.
It is not by itself a Nash theorem.  A Nash statement requires a specified
action protocol, control rights over payments, liquidity constraints,
information, side payments, and the actions available to excluded creditors.

The results suggest three possible institutional rules.

\begin{enumerate}[label=(\alph*)]
  \item \textbf{Strict standstill.}  Freeze all pre-clearing transfers.  This
  prevents preference but discards the surplus from avoiding liquidation.

  \item \textbf{Pro-rata safe harbor.}  Permit only repayments proportional to
  existing liabilities.  Theorem~\ref{thm:pro-rata} shows that this is a Pareto
  improvement for creditors in the fixed regime.

  \item \textbf{Pareto-cone safe harbor.}  Permit a targeted repayment when its
  marginal shares satisfy \eqref{eq:pareto-cone}.  The uncommitted share
  $(1-\alpha)/(1-r)$ can target a creditor near insolvency while the required
  lower-bound shares protect the others.  Implementation requires observable
  claims and an estimate of $r$.
\end{enumerate}

\section{Nonlinear liquidation losses}

A constant recovery fraction is a useful first model, but market impact should
usually depend on the amount liquidated.  Let $F(E)$ be the cash recovered from
liquidating external assets $E$, with
\begin{equation}
  0\leq F(E)\leq E.
\end{equation}
The linear model is $F(E)=\alpha E$.  Under pro-rata early repayment,
\begin{equation}
  W(t)=F(e_0-tP)+\beta I.
\end{equation}

\begin{proposition}[Pro-rata repayment with market impact]
If $F$ is differentiable, then
\begin{equation}
  \frac{dR_i}{dt}
  =L_{0i}\bigl[1-F'(e_0-tP)\bigr].
\end{equation}
Pro-rata early repayment benefits every creditor whenever the marginal
liquidation recovery satisfies $F'(E)\leq1$.
\end{proposition}

\begin{proof}
The liability share remains $L_{0i}/P$.  Therefore
\begin{equation}
  R_i(t)
  =tL_{0i}
   +\frac{L_{0i}}{P}
      \bigl[F(e_0-tP)+\beta I\bigr]+K_i.
\end{equation}
Differentiate with respect to $t$.
\end{proof}

The economically relevant primitive is thus the total recovery function $F$,
not merely a recovery fraction written as $\alpha(E)$.  If
$F(E)=\alpha(E)E$, then
\begin{equation}
  F'(E)=\alpha(E)+E\alpha'(E).
\end{equation}

\section{Discrete time}

For a dynamic version, the repayment budget must be aggregate.  At date $t$,
choose $x_i(t)\geq0$ such that
\begin{equation}
  \sum_i x_i(t)\leq e_0(t),
  \qquad
  x_i(t)\leq L_{0i}(t).
\end{equation}
The balance-sheet updates are
\begin{align}
  e_0(t+1)&=e_0(t)-\sum_i x_i(t),\\
  e_i(t+1)&=e_i(t)+x_i(t),\\
  L_{0i}(t+1)&=L_{0i}(t)-x_i(t).
\end{align}
These equations conserve external assets inside the network.  If incoming
payments to bank $0$ are unchanged, its pre-haircut solvency gap is invariant:
\begin{equation}
  \bigl[P(t+1)-e_0(t+1)-I(t+1)\bigr]
  =P(t)-e_0(t)-I(t).
\end{equation}
Thus repayment alone cannot repair bank $0$'s own fundamental deficit, though
it can alter the solvency of its creditors and hence future feedback through
the network.

A complete dynamic game would also specify the clearing time, stochastic
external-asset shocks, the timing of information, and whether actions are
observable.  Without these ingredients, claims about equilibrium timing are
not well defined.

\section{Research programme}

The calculations above lead to several concrete problems.

\subsection{A global best-response algorithm}

Theorem~\ref{thm:prepayment-jacobian} solves the smooth part of the problem.
A global algorithm should combine three types of candidates:
\begin{enumerate}[label=(\roman*)]
  \item stationary points and active faces of the strategy polytope inside a
  regular default region;
  \item solvency boundaries at which a bank changes branch in the clearing
  map;
  \item singular boundaries at which a liability row vanishes or the
  default-set resolvent ceases to exist.
\end{enumerate}
Theorem~\ref{thm:fractional-best-response} shows that the second class cannot
be omitted.  An implementable first algorithm could enumerate default sets for
small networks, use \eqref{eq:prepayment-jacobian} for local optimization, and
solve one-dimensional threshold equations along candidate edges.

\subsection{Existence of equilibrium}

The strategy space \eqref{eq:strategy-polytope} is compact and convex, but
default costs can make the utilities discontinuous at solvency boundaries.
Compactness alone therefore does not give a pure or mixed equilibrium.  Useful
next results would identify network classes for which best responses are
nonempty and convex-valued, or replace the discontinuous default rule by a
continuous recovery function and apply a fixed-point theorem.  A separate
question is whether equilibrium is restored when actions must preserve
liability proportions.

\subsection{Robust fractional optimality}

Theorem~\ref{thm:fractional-best-response} covers both recovery-parameter axes.
The next step is to characterize fractional threshold optima when
$\alpha,\beta>0$.  One expects them when the return through a newly solvent
path exceeds the direct cost of reaching the threshold, while the marginal
return beyond the threshold is below one.  A useful theorem would state these
conditions in terms of the two neighboring default-set resolvents.

\subsection{Strategic games with explicit control rights}

To turn the standstill intuition into an equilibrium theorem, specify:
\begin{enumerate}[label=(\roman*)]
  \item who proposes and authorizes a repayment;
  \item whether creditor consent is required;
  \item whether side payments are allowed;
  \item what each player observes;
  \item whether new liabilities require a matching asset transfer;
  \item how the clearing date is chosen.
\end{enumerate}
The bilateral incentive proposition then becomes an input to the game rather
than a substitute for its definition.

\subsection{Mechanism design for safe early repayment}

The local Pareto cone is the simplest targeted rule.  Further work should ask
whether its stepwise application remains Pareto improving after the effective
recovery rate and the default set change.  If selective repayment prevents a
specific contagion event, the mechanism should compare the jump in network
value at that boundary with the compensation required along the path to it.
The pro-rata rule remains the information-light benchmark because it does not
require an estimate of $r$.

\subsection{Dynamic and coalitional extensions}

With multiple action dates, a bank may fund a later prepayment from an earlier
receipt.  The game then needs explicit settlement order, information, and
clearing time.  Coalitions introduce a second problem: the debtor may have zero
equity and be indifferent while a selected creditor gains.  Side payments can
turn that weak incentive into a strict joint deviation.  A promising solution
concept is therefore coalition-proof equilibrium under the Pareto-cone safe
harbor rather than Nash equilibrium under unrestricted actions.

\section{Conclusion}

The notes point to a genuine design tension.  A complete payment freeze protects
equal treatment but can force avoidable liquidation losses.  Unrestricted
pre-clearing bargaining can preserve value but allows coalitions to redirect
recovery away from excluded creditors.

The continuous model sharpens this tension.  Inside a regular default region,
the resolvent formula makes marginal incentives computable.  Selective actions
produce rational rather than generally affine payoffs, while proportional
actions recover an affine structure.  At region boundaries, fractional
payments can be uniquely optimal because the least payment that rescues a bank
activates an entire return path through the network.

These results support a middle course for a clearing protocol: enforce
balance-sheet consistency, use pro-rata repayment when only local information
is available, and permit targeted departures when the Pareto-cone constraints
can be verified.  The main mathematical problem left is global: combine smooth
optimization inside default regions with the discrete solvency transitions
between them, then establish equilibrium under the resulting discontinuous
payoffs.

\begin{thebibliography}{99}

\bibitem{EisenbergNoe2001}
L.~Eisenberg and T.~H. Noe.
\newblock Systemic risk in financial systems.
\newblock \emph{Management Science}, 47(2):236--249, 2001.

\bibitem{RogersVeraart2013}
L.~C.~G. Rogers and L.~A.~M. Veraart.
\newblock Failure and rescue in an interbank network.
\newblock \emph{Management Science}, 59(4):882--898, 2013.

\bibitem{AminiEtAl2016}
H.~Amini, D.~Filipovi\'c, and A.~Minca.
\newblock To fully net or not to net: Adverse effects of partial multilateral
netting.
\newblock \emph{Operations Research}, 64(5):1135--1142, 2016.

\bibitem{SchuldenzuckerSeuken2020}
S.~Schuldenzucker and S.~Seuken.
\newblock Portfolio compression in financial networks: Incentives and
systemic risk.
\newblock In \emph{Proceedings of the 21st ACM Conference on Economics and
Computation}, page 79, 2020.

\bibitem{BertschingerEtAl2020}
N.~Bertschinger, M.~Hoefer, and D.~Schmand.
\newblock Strategic payments in financial networks.
\newblock In \emph{11th Innovations in Theoretical Computer Science
Conference}, volume 151 of \emph{LIPIcs}, pages 46:1--46:16, 2020.

\bibitem{PappWattenhofer2020}
P.~A. Papp and R.~Wattenhofer.
\newblock Network-aware strategies in financial systems.
\newblock In \emph{47th International Colloquium on Automata, Languages, and
Programming}, volume 168 of \emph{LIPIcs}, pages 91:1--91:17, 2020.

\bibitem{KusnetsovVeraart2019}
M.~Kusnetsov and L.~A.~M. Veraart.
\newblock Interbank clearing in financial networks with multiple maturities.
\newblock \emph{SIAM Journal on Financial Mathematics}, 10(1):37--67, 2019.

\bibitem{LiuStaum2010}
M.~Liu and J.~Staum.
\newblock Sensitivity analysis of the Eisenberg--Noe model of contagion.
\newblock \emph{Operations Research Letters}, 38(5):489--491, 2010.

\bibitem{FeinsteinEtAl2018}
Z.~Feinstein, W.~Pang, B.~Rudloff, E.~Schaanning, S.~Sturm, and M.~Wildman.
\newblock Sensitivity of the Eisenberg--Noe clearing vector to individual
interbank liabilities.
\newblock \emph{SIAM Journal on Financial Mathematics}, 9(4):1286--1325,
2018.

\bibitem{ZhouEtAl2024}
H.~Zhou, Y.~Wang, K.~Varsos, N.~Bishop, R.~Savani, A.~Calinescu, and
M.~Wooldridge.
\newblock A strategic analysis of prepayments in financial credit networks.
\newblock In \emph{Proceedings of the Thirty-Third International Joint
Conference on Artificial Intelligence}, pages 3040--3047, 2024.

\end{thebibliography}

\end{document}
